\subsection{模的对称代数的定义}

定义 1. 设 A 为交换环，M 为 A-模。M 的对称代数，记为 S(M) 或 Sym(M)，或 \(S_{A}(M)\)，是张量代数 T(M) 关于双边理想 \(\mathfrak{J}'\) （也记作 \(\mathfrak{J}_{\mathbf{M}}'\) ）的商代数，该理想由 T(M) 的元素 \(xy - yx = x \otimes y - y \otimes x\) 生成，其中 \(x\) 和 \(y\) 遍历 M。

由于理想 \(\mathfrak{J}'\) 由次数为 2 的齐次元素生成，它是一个分次理想（II，§11，第3号，命题2）；我们记 \(\mathfrak{J}_n' = \mathfrak{J}' \cap T^n(M)\) ；代数 \(S(M)\) 于是由诸 \(S^n(M) = T^n(M)/\mathfrak{J}_n'\) 组成的分次（称为典范分次）进行分次。现在 \(\mathfrak{J}_0' = \mathfrak{J}_1' = \{0\}\) ，因此 \(S^0(M)\) 与 A 典范等同，而 \(S^1(M)\) 与 \(T^1(M) = M\) 等同；在后续内容中我们始终进行这些等同，并用 \(\phi'\) 或 \(\phi_M'\) 表示典范单射 \(M \to S(M)\) 。

命题1. 代数 \(\mathbf{S}(\mathbf{M})\) 是交换的。

根据定义 \(\phi'(x)\phi'(y) = \phi'(y)\phi'(x)\) （对 \(x, y\) 属于 \(\mathbf{M}\) ）；并且由于元素 \(\phi'(x)\) （其中 \(x\) 遍历 \(\mathbf{M}\) ）生成 \(\mathbf{S}(\mathbf{M})\) ，结论由 §1 第 7 号得出。

命题2. 设E是一个A-代数，且f: M → E是一个A-线性映射，使得

\[
f (x) f (y) = f (y) f (x) \text {   for   all   } x, y \text {   in   } \mathbf {M}.\tag{1}
\]

存在唯一的 A-代数同态 \(g \colon \mathbb{S}(\mathbf{M}) \to \mathbf{E}\) ，使得 \(f = g \circ \phi'\) 。

换言之， \((\mathbb{S}(\mathbf{M}), \phi')\) 是泛映射问题（集合论，IV，§3，第1号）的一个解，其中 \(\Sigma\) 是A-代数结构的类， \(\alpha\) -映射是满足(1)的从A-模M到A-代数的线性映射。

\(g\) 的唯一性由 \(\phi'(\mathbf{M}) = \mathbf{M}\) 生成 \(\mathbf{S}(\mathbf{M})\) 这一事实推出。为证明 \(g\) 的存在性，注意根据 §5 第 1 号命题 1，存在一个 A-代数同态 \(g_1: \mathbf{T}(\mathbf{M}) \to \mathbf{E}\) ，使得 \(f = g_1 \circ \phi\) ；

只需证明 \(g_1\) 在理想 \(\mathfrak{J}'\) 上为零即可，因为此时，若 \(p: \mathsf{T}(\mathbf{M}) \to \mathsf{S}(\mathbf{M}) = \mathsf{T}(\mathbf{M}) / \mathfrak{J}'\) 是典范同态，则可写 \(g_1 = g \circ p\) ，其中 \(g: \mathsf{S}(\mathbf{M}) \to \mathbf{E}\) 是一个代数同态，而结论将由 \(p \circ \phi = \phi'\) 这一事实得出。现在，\(g_1\) 的核是一个双边理想，由（1）及关系 \(g_1 \circ \phi = f\) ，它包含元素 \(x \otimes y - y \otimes x\) （对 \(x, y\) 属于 \(\mathbf{M}\) ）。这就完成了证明。

注：(1) 设 E 是一个 Z 型分次 A-代数，其分次为 \((\mathrm{E}_{n})\) ，并设线性映射 f （假定满足 (1)）满足

\[
f (\mathbf {M}) \subset \mathrm{E} _ {1}.\tag{2}
\]

则关系式 \(g(x_1x_2 \ldots x_p) = f(x_1)f(x_2) \ldots f(x_p)\) （其中 \(x_i \in \mathbf{M}\) ）表明 \(g(\mathbf{S}^p(\mathbf{M})) \subset \mathrm{E}_p\) 对所有 \(p \geqslant 0\) 成立，因此 \(g\) 是分次代数同态。

\begin{enumerate}
\def\labelenumi{(\arabic{enumi})}
\setcounter{enumi}{1}
\item
  \(\mathsf{S}(\mathbf{M})\) 的每一个元素都是形如 \(x_{1}x_{2}\ldots x_{n}\) （其中 \(x_{i}\) 属于 \(\mathbf{M}\) ）的乘积之和；应注意不要把 \(\mathsf{S}(\mathbf{M})\) 中所取的此类乘积与 \(\mathsf{T}(\mathbf{M})\) 中所取的类似乘积相混淆。
\item
  若 \(n!\) .1 在 A 中可逆，A-模 \(S^n(M)\) 由形如 \(x^n\) （其中 \(x \in M\) ）的元素生成；这由上述注记及 I，§8，第 2 号，命题 2 得出。
\end{enumerate}

